\documentclass[11pt,a4paper]{article}
\usepackage[utf8]{inputenc}
\usepackage[T1]{fontenc}
\usepackage{lmodern}
\usepackage[french]{babel}
\usepackage[margin=2.3cm]{geometry}
\usepackage{amsmath,amssymb,mathtools}
\usepackage{bm}
\usepackage{enumitem}
\usepackage[dvipsnames]{xcolor}
\usepackage[most]{tcolorbox}
\usepackage{fancyhdr}
\pagestyle{fancy}\fancyhf{}
\lhead{\small Chapitre 7 -- La non-sphéricité des erreurs}\rhead{\small MCG plus efficace que MCO}\cfoot{\thepage}
\setlength{\headheight}{14pt}
\setlength{\parskip}{0.5em}\setlength{\parindent}{0pt}

\newcommand{\vy}{\mathbf{y}}\newcommand{\vX}{\mathbf{X}}\newcommand{\vb}{\mathbf{b}}\newcommand{\vc}{\mathbf{c}}
\newcommand{\Xs}{\mathbf{X}^{*}}\newcommand{\W}{\mathbf{W}}
\newcommand{\veps}{\bm{\epsilon}}\newcommand{\vbeta}{\bm{\beta}}\newcommand{\tbeta}{\tilde{\bm{\beta}}}
\newcommand{\I}{\mathbf{I}}\newcommand{\Pm}{\mathbf{P}}\newcommand{\Om}{\bm{\Omega}}
\newcommand{\A}{\mathbf{A}}\newcommand{\B}{\mathbf{B}}\newcommand{\D}{\mathbf{D}}
\newcommand{\Mw}{\mathbf{M}_W}
\newcommand{\Var}{\operatorname{Var}}\newcommand{\E}{\mathbb{E}}

\newtcolorbox{idee}[1][]{colback=Cerulean!6,colframe=Cerulean!70!black,fonttitle=\bfseries,title={Idée de la démonstration},#1}
\newtcolorbox{resultat}[1][]{colback=ForestGreen!6,colframe=ForestGreen!70!black,fonttitle=\bfseries,title={Résultat de l'étape},#1}
\newtcolorbox{intuition}[1][]{colback=Goldenrod!8,colframe=Goldenrod!70!black,fonttitle=\bfseries,title={Intuition},#1}
\newtcolorbox{attention}[1][]{colback=BrickRed!5,colframe=BrickRed!75!black,fonttitle=\bfseries,title={Attention},#1}

\begin{document}

\begin{center}
{\LARGE\bfseries Les MCG sont plus efficaces que les MCO}\\[6pt]
{\large Démonstration pas à pas (chapitre 7, section 3.1 : « laissée comme exercice »)}
\end{center}

\section*{Le point de départ (cours, section 3.1)}

Modèle : $\vy=\vX\vbeta+\veps$, \ $\E[\veps\mid\vX]=\mathbf 0$, \ $\E[\veps\veps'\mid\vX]=\sigma^2\Om$, avec $\Om$ symétrique et définie positive.

\[
\Var(\vb\mid\vX)=\sigma^2(\vX'\vX)^{-1}(\vX'\Om\vX)(\vX'\vX)^{-1}\quad\text{(MCO)},\qquad
\Var(\tbeta\mid\vX)=\sigma^2(\vX'\Om^{-1}\vX)^{-1}\quad\text{(MCG)}.
\]
Il faut montrer que $\Var(\vb\mid\vX)-\Var(\tbeta\mid\vX)$ est semi-définie positive. Le cours propose de passer par les inverses :
\[
\Var(\tbeta\mid\vX)^{-1}-\Var(\vb\mid\vX)^{-1}=\frac1{\sigma^2}\,\D,\qquad
\D=\vX'\Om^{-1}\vX-(\vX'\vX)(\vX'\Om\vX)^{-1}(\vX'\vX).
\]
L'inverse de $\Var(\vb)$ s'obtient avec $(\A\B\mathbf C)^{-1}=\mathbf C^{-1}\B^{-1}\A^{-1}$ :
\[
\Var(\vb\mid\vX)^{-1}=\frac1{\sigma^2}(\vX'\vX)(\vX'\Om\vX)^{-1}(\vX'\vX).
\]

\begin{idee}
Réécrire $\D$ sous la forme $\Xs{}'\,\Mw\,\Xs$, où $\Mw$ est une \textbf{matrice de projection} (symétrique, idempotente) du type $\mathbf M$ du chapitre 2. Une telle forme est toujours semi-définie positive, car elle s'écrit comme un « carré » : $(\Mw\Xs)'(\Mw\Xs)$. C'est la même astuce que dans la preuve de Gauss-Markov.
\end{idee}

\section*{Étape 1 : la matrice $\Pm$ du cours}

$\Om$ étant symétrique et définie positive, il existe une matrice $\Pm$ ($n\times n$) inversible telle que (cours, éq. (6)--(8)) :
\[
\Pm'\Pm=\Om^{-1}\qquad\text{et}\qquad \Om=(\Pm'\Pm)^{-1}=\Pm^{-1}(\Pm')^{-1}.
\]

\section*{Étape 2 : deux matrices transformées}

On définit
\[
\Xs=\Pm\vX\qquad\text{(les régresseurs transformés « MCG »)},\qquad\qquad \W=(\Pm')^{-1}\vX .
\]
Ce sont deux matrices $n\times K$.

\section*{Étape 3 : réécrire les trois blocs de $\D$}

\textbf{Le premier terme.}
\[
\vX'\Om^{-1}\vX=\vX'\Pm'\Pm\vX=(\Pm\vX)'(\Pm\vX)=\Xs{}'\Xs .
\]

\textbf{Le terme du milieu.} Comme $(\Pm^{-1})'=(\Pm')^{-1}$, on a $\vX'\Pm^{-1}=\big((\Pm')^{-1}\vX\big)'=\W'$, donc
\[
\vX'\Om\vX=\vX'\Pm^{-1}(\Pm')^{-1}\vX=\big((\Pm')^{-1}\vX\big)'\big((\Pm')^{-1}\vX\big)=\W'\W .
\]

\textbf{Le terme $\vX'\vX$.} On insère $\Pm'(\Pm')^{-1}=\I_n$ au milieu :
\[
\vX'\vX=\vX'\,\Pm'(\Pm')^{-1}\,\vX=(\Pm\vX)'\big((\Pm')^{-1}\vX\big)=\Xs{}'\W,
\qquad\text{et en transposant :}\qquad \vX'\vX=\W'\Xs .
\]

\begin{resultat}
$\vX'\Om^{-1}\vX=\Xs{}'\Xs$, \qquad $\vX'\Om\vX=\W'\W$, \qquad $\vX'\vX=\Xs{}'\W=\W'\Xs$.
\end{resultat}

\section*{Étape 4 : faire apparaître une matrice de projection}

On remplace dans $\D$ :
\begin{align*}
\D&=\Xs{}'\Xs-\Xs{}'\W\,(\W'\W)^{-1}\,\W'\Xs\\
&=\Xs{}'\Big[\I_n-\W(\W'\W)^{-1}\W'\Big]\Xs\\
&=\Xs{}'\,\Mw\,\Xs ,
\end{align*}
où $\Mw=\I_n-\W(\W'\W)^{-1}\W'$ est la matrice qui crée les résidus d'une régression sur $\W$ (la matrice $\mathbf M$ du chapitre 2, avec $\W$ à la place de $\vX$). ($\W'\W=\vX'\Om\vX$ est bien inversible : $\Om$ est définie positive et $\vX$ de plein rang.)

$\Mw$ est \textbf{symétrique} ($\Mw'=\Mw$) et \textbf{idempotente} ($\Mw\Mw=\Mw$). Donc
\[
\D=\Xs{}'\Mw\Mw\Xs=\Xs{}'\Mw'\Mw\Xs=(\Mw\Xs)'(\Mw\Xs).
\]
Pour tout vecteur $\vc\in\mathbb R^K$ :
\[
\vc'\D\vc=(\Mw\Xs\vc)'(\Mw\Xs\vc)=\big\|\Mw\Xs\vc\big\|^2\ \ge\ 0 .
\]

\begin{resultat}
$\D=(\Mw\Xs)'(\Mw\Xs)$ est \textbf{semi-définie positive}. Donc $\Var(\tbeta\mid\vX)^{-1}-\Var(\vb\mid\vX)^{-1}=\D/\sigma^2$ est semi-définie positive.
\end{resultat}

\section*{Étape 5 : passer des inverses aux variances}

On utilise un résultat d'algèbre matricielle (Greene, annexe A) :

\begin{tcolorbox}[colback=white,colframe=gray!60!black]
Si $\A$ et $\B$ sont symétriques définies positives, alors
\[
\A-\B\ \text{semi-définie positive}\quad\Longleftrightarrow\quad \B^{-1}-\A^{-1}\ \text{semi-définie positive}.
\]
C'est l'analogue matriciel de : si $a\ge b>0$, alors $\frac1b\ge\frac1a$.
\end{tcolorbox}

Avec $\A=\Var(\tbeta\mid\vX)^{-1}$ et $\B=\Var(\vb\mid\vX)^{-1}$ : on vient de montrer que $\A-\B$ est semi-définie positive, donc $\B^{-1}-\A^{-1}=\Var(\vb\mid\vX)-\Var(\tbeta\mid\vX)$ l'est aussi.

\begin{resultat}[title={Conclusion}]
\[
\Var(\vb\mid\vX)-\Var(\tbeta\mid\vX)\quad\text{est semi-définie positive.}
\]
Chaque coefficient, et plus généralement toute combinaison linéaire $\vc'\vbeta$, est estimé avec une variance au moins aussi petite par les MCG que par les MCO : $\Var(\vc'\tbeta)\le\Var(\vc'\vb)$. \textbf{Les MCG sont plus efficaces que les MCO.} \hfill$\blacksquare$
\end{resultat}

\section*{Remarques}

\begin{attention}[title={« Définie » ou « semi-définie » positive ?}]
Le cours écrit « positive définie », mais en général la différence est seulement \textbf{semi-définie} positive. Elle est nulle quand $\Mw\Xs=\mathbf 0$, et alors MCO et MCG coïncident. C'est le cas :
\begin{itemize}[nosep]
\item si $\Om=\I_n$ (erreurs sphériques) ;
\item dans l'exercice 1 du chapitre, si $\sigma_1^2=\sigma_2^2$ ;
\item dans l'exercice 2 avec deux observations : les MCG donnent la moyenne simple, comme les MCO.
\end{itemize}
\end{attention}

\textbf{Preuve plus courte, par Gauss-Markov.} Dans le modèle transformé $\Pm\vy=\Pm\vX\vbeta+\Pm\veps$, les erreurs sont sphériques : $\E[\Pm\veps\veps'\Pm'\mid\vX]=\sigma^2\I_n$ (cours, éq. (9)). Le théorème de Gauss-Markov s'applique donc à ce modèle, et son estimateur MCO est $\tbeta$ (éq. (10)). Or l'estimateur des MCO du modèle d'origine s'écrit
\[
\vb=(\vX'\vX)^{-1}\vX'\vy=(\vX'\vX)^{-1}\vX'\Pm^{-1}\,(\Pm\vy) :
\]
il est \emph{linéaire} en $\Pm\vy$ et \emph{sans biais}. Par Gauss-Markov, $\Var(\vb\mid\vX)\ge\Var(\tbeta\mid\vX)$. La démonstration en 5 étapes ci-dessus fait ce calcul explicitement, sans invoquer le théorème.

\begin{intuition}
Les MCO donnent le même poids à toutes les observations, même aux plus bruitées (hétéroscédasticité) ou aux plus redondantes (autocorrélation : deux observations corrélées apportent moins d'information que deux observations indépendantes). Les MCG pondèrent chaque observation par l'information qu'elle apporte réellement, à travers $\Om^{-1}$. Exemple de l'exercice 1 avec $\sigma_1^2=1$ et $\sigma_2^2=9$ : poids $(0{,}9\ ;\ 0{,}1)$ pour les MCG au lieu de $(0{,}5\ ;\ 0{,}5)$, et une variance de 0,9 au lieu de 2,5.
\end{intuition}

\end{document}
